\documentclass[12pt]{article}
\usepackage[utf8]{inputenc}
\usepackage[T1]{fontenc}
\usepackage{hyperref}
\title{Wu Chinese Language Resources for K-12 Curriculum Development}
\author{Compiled by Claude Sonnet 5.5 (claude-sonnet-5-5) for Andr\'as Kornai}
\date{}
\begin{document}
\maketitle

\section{Introduction}
Wu Chinese (ISO 639-3 wuu, Glottocode wuch1236) is spoken mainly in Shanghai, Zhejiang and southern Jiangsu. Omniglot gives about 90 million speakers; the survey figure is about 77 million. Suzhou was traditionally the prestige variety, and Shanghainese has since overtaken it. See \url{https://www.omniglot.com/chinese/wu.htm}. Wu is a family of many varieties, so a curriculum must choose a variety or varieties.

\section{Primary Educational Resources}
Searches found no school textbooks in Wu. They found learner-oriented material. The Wikibooks textbook \emph{Wu Chinese} teaches the Suzhou dialect, spoken and written: \url{https://en.wikibooks.org/wiki/Wu_Chinese/Printable_version}. Other learner courses were mentioned in search summaries (ZanHei, ChinesePod Shanghainese), but no URLs for them were seen, so they are not linked here.

\section{Linguistic and Reference Materials}
Hippocrene's \emph{Shanghainese-English/English-Shanghainese Dictionary and Phrasebook} has about 4,000 entries plus brief grammar and pronunciation sections: \url{https://alc.rutgers.edu/people/faculty-bookshelf/bookshelf-list/77-shanghainese-englishenglish-shanghainese-dictionary-and-phrasebook}. Lexilogos lists Wu dictionaries and links, including Creamer's \emph{Chinese--English Dictionary of the Wu Dialect} (1991): \url{https://www.lexilogos.com/english/wu_dictionary.htm} and \url{https://www.lexilogos.com/wu_dictionnaire.htm}. A 2006 forum thread notes that Wu dictionaries use IPA because no standard romanization exists; this may be out of date: \url{https://www.chinese-forums.com/forums/topic/7457-a-dictionary-of-wu/}.

\section{Religious and Community Resources}
No scripture or community organisation was identified in these searches.

\section{Digital and Audio Resources}
A Reddit-hosted wiki on Shanghainese was found: \url{https://r.datuan.dev/r/shanghainese/wiki/index}. The Wugniu online dictionary (Suzhou, Wuxi, Hangzhou, Ningbo and others) was mentioned in search summaries, but no URL was seen.

\section{Curriculum Development Context}
Wu is not known to be taught in schools. Existing materials are aimed at adult learners and linguists, not K-12. A curriculum would need a variety-specific orthography decision, with Mandarin as the language of school instruction.

\end{document}